\documentclass[11pt]{article}
% Shared design for the two minimal papers.
% It follows the layout of the guide "The Riemann Hypothesis and its generalizations"
% (KNT with Fable, 9 November 2026): a designed title block, a short story list at the start,
% coloured headings, boxed statements, idea boxes, and annotated figures.
%
% A paper sets \papertitle, \papersubtitle and \paperrunning before \begin{document}
% and calls \maketitleblock after it.

\usepackage[letterpaper,left=1.12in,right=1.12in,top=1.05in,bottom=1.1in]{geometry}
\usepackage{amsmath,mathtools,amssymb,amsthm}
\usepackage{newpxtext}
\usepackage[varbb]{newpxmath}
\usepackage{fontspec}
\setsansfont{texgyreheros}[Extension=.otf, UprightFont=*-regular, BoldFont=*-bold,
  ItalicFont=*-italic, BoldItalicFont=*-bolditalic, Scale=0.88]
\usepackage{microtype}
\usepackage{parskip}
\usepackage{xcolor}
\usepackage{booktabs,array}
\usepackage{enumitem}
\usepackage{multicol}
\usepackage[titles]{tocloft}
\usepackage{tikz}
\usetikzlibrary{decorations.markings,arrows.meta,positioning,calc}
\usepackage[most]{tcolorbox}
\usepackage{fancyhdr}
\usepackage{titlesec}
\usepackage[font=small,labelfont={bf,sf},labelsep=period,skip=6pt]{caption}
\usepackage{url}

\definecolor{accent}{HTML}{1B4F72}
\definecolor{warm}{HTML}{C0562A}
\definecolor{slate}{HTML}{56636E}
\definecolor{ink}{HTML}{1E1E1E}

\usepackage[colorlinks=true,linkcolor=accent,citecolor=accent,urlcolor=accent]{hyperref}

% Lean identifiers: typewriter, breakable at dots and underscores.
\DeclareUrlCommand\lean{\urlstyle{tt}}

% ---------- layout ----------
\linespread{1.06}
\setlength{\parskip}{0.55em plus 0.1em minus 0.05em}
\setlist{itemsep=0.25em, topsep=0.2em, parsep=0.1em}
\titleformat{\section}{\sffamily\bfseries\Large\color{accent}}{\thesection}{0.7em}{}
\titleformat{\subsection}{\sffamily\bfseries\normalsize\color{ink}}{\thesubsection}{0.6em}{}
\titleformat{\paragraph}[runin]{\sffamily\bfseries\small\color{accent}}{}{0pt}{}[\enspace]
\titlespacing*{\section}{0pt}{22pt}{5pt}
\titlespacing*{\subsection}{0pt}{13pt}{1pt}
\titlespacing*{\paragraph}{0pt}{9pt}{0pt}
\pagestyle{fancy}
\fancyhf{}
\renewcommand{\headrulewidth}{0pt}
\fancyhead[L]{\sffamily\footnotesize\color{slate}\paperrunning}
\fancyhead[R]{\sffamily\footnotesize\color{slate}\thepage}
\setlength{\headheight}{14pt}
\setcounter{tocdepth}{1}
\setlength{\cftbeforesecskip}{2pt}
\renewcommand{\cftsecfont}{\small}
\renewcommand{\cftsecpagefont}{\small\color{slate}}
\renewcommand{\cftsecleader}{\cftdotfill{\cftdotsep}}
\setlength{\cftsecnumwidth}{1.7em}
\setlength{\emergencystretch}{2em}
\numberwithin{equation}{section}

% ---------- title block ----------
\newcommand{\papertitle}{}
\newcommand{\papersubtitle}{}
\newcommand{\paperrunning}{}
\newcommand{\paperauthors}{KNT\enspace\textperiodcentered\enspace Astra\enspace\textperiodcentered\enspace Fable~5.2}
\newcommand{\paperdate}{Draft of 11 November 2026}
\newcommand{\maketitleblock}{%
  \thispagestyle{plain}%
  {\sffamily\bfseries\color{accent}\fontsize{25}{29}\selectfont\raggedright\papertitle\par}%
  \vspace{5pt}%
  {\sffamily\Large\color{slate}\raggedright\papersubtitle\par}%
  \vspace{8pt}%
  {\color{accent}\hrule height 1.2pt}%
  \vspace{7pt}%
  {\sffamily\small\color{ink}\paperauthors\hfill\color{slate}\paperdate\par}%
  \vspace{10pt}}

% Small headings inside the front matter.
\newcommand{\fronthead}[1]{\par\vspace{6pt}{\sffamily\bfseries\normalsize\color{ink}#1}\par\vspace{-2pt}}
\makeatletter
\newcommand{\contentsblock}{%
  \par\noindent\begin{minipage}{\linewidth}%
  \fronthead{Contents}%
  \setlength{\parskip}{0pt}\setlength{\columnsep}{2.2em}\raggedright
  \begin{multicols}{2}
  \@starttoc{toc}
  \end{multicols}
  \end{minipage}\par}
\makeatother

% ---------- statements ----------
% Theorems, lemmas, propositions and corollaries share one counter.
% Each statement is in a tinted box with an accent rule at the left.
\newtheoremstyle{boxed}{0pt}{0pt}{\normalfont}{0pt}{\sffamily\bfseries\small\color{accent}}{.}{0.6em}{}
\theoremstyle{boxed}
\newtheorem{theorem}{Theorem}[section]
\newtheorem{lemma}[theorem]{Lemma}
\newtheorem{proposition}[theorem]{Proposition}
\newtheorem{corollary}[theorem]{Corollary}

\tcbset{guidebox/.style={enhanced, parbox=false, boxrule=0pt, frame hidden,
  arc=1.5pt, left=11pt, right=10pt, top=7pt, bottom=7pt, before skip=11pt, after skip=9pt}}
% Not breakable: a label at the start of a box that breaks at the top of a page can be lost.
\tcbset{statementbox/.style={guidebox, colback=accent!6,
  borderline west={3pt}{0pt}{accent}}}
\tcolorboxenvironment{theorem}{statementbox}
\tcolorboxenvironment{lemma}{statementbox}
\tcolorboxenvironment{proposition}{statementbox}
\tcolorboxenvironment{corollary}{statementbox}

% The main result of a paper: a framed box with a title bar.
% \begin{maintheorem}{name} uses the counter of the theorems, and \label works after it.
\newtcolorbox{mainresult}[1]{enhanced, parbox=false, colback=white, colframe=accent, boxrule=0.9pt,
  arc=1.5pt, left=11pt, right=11pt, top=8pt, bottom=8pt, title={#1},
  fonttitle=\sffamily\bfseries\small, colbacktitle=accent, coltitle=white,
  before skip=12pt, after skip=12pt}
\newenvironment{maintheorem}[1]%
  {\refstepcounter{theorem}\begin{mainresult}{Theorem~\thetheorem\ (#1)}}%
  {\end{mainresult}}

% ---------- proofs ----------
% A proof has a thin rule at the left, so that the eye finds the statements.
\renewcommand{\proofname}{\upshape\sffamily\bfseries\small\color{slate}Proof}
\renewcommand{\qedsymbol}{\textcolor{slate}{\rule{0.55em}{0.55em}}}
\tcolorboxenvironment{proof}{enhanced, breakable, parbox=false, boxrule=0pt, frame hidden,
  colback=white, borderline west={1.2pt}{0pt}{slate!35}, left=11pt, right=0pt, top=1pt, bottom=2pt,
  before skip=6pt, after skip=11pt}
% A step or a case inside a proof.
\newcommand{\step}[1]{{\sffamily\bfseries\small\color{slate}#1}}

% ---------- idea, notation and fine-print boxes ----------
\newtcolorbox{idea}[1]{guidebox, breakable, colback=warm!6, borderline west={3pt}{0pt}{warm!80},
  before upper={{\sffamily\bfseries\footnotesize\color{warm!75!black}IDEA\enspace\textbar\enspace}{\sffamily\bfseries\small #1}\par}}
\newtcolorbox{notation}[1]{guidebox, breakable, colback=black!4, borderline west={3pt}{0pt}{slate!60},
  fontupper=\small,
  before upper={{\sffamily\bfseries\footnotesize\color{slate}NOTATION\enspace\textbar\enspace}{\sffamily\bfseries\small #1}\par}}
\newtcolorbox{fineprint}{guidebox, breakable, colback=white, borderline west={1.5pt}{0pt}{black!25},
  fontupper=\footnotesize\color{black!78}, left=9pt, top=1pt, bottom=1pt,
  before upper={{\sffamily\bfseries Fine print.}\enspace}}

% ---------- figures ----------
\tikzset{
  vtx/.style={circle, fill=ink, inner sep=1.35pt},
  port/.style={circle, draw=accent, fill=white, line width=0.8pt, inner sep=1.5pt},
  edge/.style={line width=0.55pt, draw=ink},
  hot/.style={line width=1.1pt, draw=warm},
  cool/.style={line width=1.1pt, draw=accent},
  region/.style={draw=accent, fill=accent!7, line width=0.6pt, rounded corners=6pt},
  note/.style={font=\sffamily\footnotesize, text=slate},
}

% ---------- macros of the mathematics ----------
\newcommand{\Z}{\mathcal{Z}}
\newcommand{\PP}{\mathbb{P}}
\newcommand{\EE}{\mathbb{E}}
\newcommand{\comp}{\operatorname{comp}}
\newcommand{\tr}{\operatorname{tr}}
\newcommand{\lgk}[1]{\log^{(#1)}}
\newcommand{\lgs}{\log^{*}}
\newcommand{\sm}{\setminus}
\newcommand{\half}{\tfrac12}
\newcommand{\cO}{\mathcal{O}}


\renewcommand{\papertitle}{An audit of Astra's notes,\\ and one combined form of the proof}
\renewcommand{\papersubtitle}{A fourth research note on ``The maximum number of cycles in a graph of large rank''}
\renewcommand{\paperrunning}{An audit of Astra's notes}
\renewcommand{\paperauthors}{Fable~5.2, for KNT}
\renewcommand{\paperdate}{Research note of 12 November 2026}
\hypersetup{pdftitle={An audit of Astra's notes, and one combined form of the proof},pdfauthor={Fable 5.2, for KNT}}

\newcommand{\M}[1]{[M,~#1]}
\newcommand{\T}[1]{[T,~#1]}

\begin{document}
\maketitleblock

KNT gave me notes of Astra with eight proposals for a shorter proof, and asked which of them can
be incorporated. This note has the audit, and it describes one copy of the paper that contains
the sound parts together with the reductions of my three notes [N1], [N2], [N3]. [M] is the
draft of 42~pages; Astra's notes use its numbers. [T] is the new copy,
\texttt{trial-combined.pdf}. The manuscript is not changed.

\fronthead{The result in eight sentences}
\begin{enumerate}[leftmargin=1.6em]
\item \emph{The audit.} I checked every step of the eight proposals, and I found no error. One
step needs an exact reading, which Section~\ref{sec:path} gives.
\item \emph{Agreement.} Three proposals are the same as reductions of my notes: the marked set,
the gap theorem as a corollary, and the recursion from a length $d$ to a length of order
$\log d$. Astra's notes do not mention [N1], [N2], [N3], so I take this as two independent
derivations.
\item \emph{New and in the copy.} The lemma on a shortest path, which needs no hypothesis on the
number of paths; the proof of rigidity by a limit; the cut of size three that excludes three
full pairs at a port; the upper bound for pairs with two links; a shorter proof of nesting; a
shorter window step in the lemma on layers.
\item \emph{Sound, but not in the copy.} Caps for cuts of size two and three: measured, 11~lines
more. Connected blocks in the chain: measured, 1~line more. The form of the theorem on cuts of
size four with two sides: not necessary when one side is a cycle of length four. The recursion
with two deleted sets: the form of [N1] has one family of cycles.
\item \emph{Measured.} The copy [T] has 41~pages, 1577 typeset lines, 39 numbered statements and
three cases. The manuscript has 42, 1624, 40 and five. The copy of [N3] has 41, 1587, 39 and
three.
\item \emph{A price of [N3] is gone.} With the limit proof, the gap constant is used at one
place. Proposition~6.3 keeps the proof and the numbers of the manuscript, and the lemma on a
full star is not necessary.
\item \emph{Astra's notes keep the case of bounded girth.} With [N2] the statements that the
notes list as retained (\M{Lemma~5.5}, \M{Corollary~5.6}, \M{Section~7}) are not necessary.
The copy [T] uses [N2].
\item \emph{Status.} All new proofs are on paper. Nothing of this note or of Astra's notes has a
formal proof, and the text of [T] differs from the Lean proof in more places than [M].
\end{enumerate}

\begin{notation}{the status of this note}
\begin{itemize}[leftmargin=1.2em]
\item \emph{Checked on paper.} Each proposal, in the form of Astra's notes. The copy [T] has the
proofs of the adopted parts in full, and I read the typeset text of its Sections~4 to~7 again.
\item \emph{Tested by computer.} The script \texttt{checks/astra\_audit\_check.py}
(Section~\ref{sec:tests}). The tests cover small graphs and the arithmetic. They do not replace
a proof.
\item \emph{Measured.} Every number of lines in this note is from a built copy.
\item \emph{Not done.} No formal proof. No reader other than Astra, for its own proposals, and
me. I did not build a copy without [N2] (Section~\ref{sec:open}).
\end{itemize}
\end{notation}

\contentsblock

\section{Verdicts}\label{sec:verdicts}

\begin{center}\small
\begin{tabular}{@{}p{1.1cm}p{4.6cm}p{1.0cm}p{2.1cm}p{5.0cm}@{}}
\toprule
 & \textsf{\textbf{Proposal}} & \textsf{\textbf{Sound}} & \textsf{\textbf{In [T]}} & \textsf{\textbf{Measured effect, remark}}\\
\midrule
1 & Marked set; no edge recursion; five marked vertices at the end & yes & yes & same as reduction~4 of [N3]; $-2$ lines\\
2 & Window step with the threshold $m^{3/4}$ & yes & yes & Section~4: $-5$ lines\\
3 & A shortest path: counts without a hypothesis; cycles by their edges on the path & yes & yes & Section~5: $-1$ line; no ordinary layers, no matchings, one conclusion in \M{Proposition~5.3}\\
4 & A gap with some fixed constant, without induction & yes & yes & same as reduction~1 of [N3]; $-24$ lines there\\
5 & Caps for cuts of size two and three; the bound $2^{r-1}+420\cdot2^{r/2}$ & yes & no (option) & $+11$ lines\\
6 & Rigidity by limits; three full pairs give a cut of size three & yes & yes, for a cycle of length four & Section~6: $-6$ lines; no lemma on a full star\\
7.1 & Shorter proof of nesting & yes & yes & with 7.2: $-3$ lines\\
7.1 & Connected blocks & yes & no (option) & $+1$ line\\
7.2 & Two links as an upper bound for $W_2$ & yes & yes & shorter statement; shorter proof for heavy families\\
7.3--7.5 & Packing at logarithmic girth; recursion $d\mapsto C_0\log(Qd+2)$; direct end & yes & no & same recursion as [N1]; [N1] has one family of cycles, and [N2] needs another base\\
8.1 & Coarse bound for the dyadic scales & yes & no & no dyadic scales in [N1]\\
8.2 & Fixed blocks in \M{Proposition~5.3} & yes & no & Astra does not recommend it\\
\bottomrule
\end{tabular}
\end{center}

\noindent The effects are differences of typeset lines against the copy of [N3] (lean form
of~[N2] with the four reductions), section by section. The introduction of [T] is five lines
longer, because it states the origin of the new parts.

\section{The audit, proposal by proposal}\label{sec:audit}

\subsection{Proposal 1: the marked set}

For a connected cubic multigraph without a bridge and a set $P$ of $k$ vertices, Astra proves
$\beta(G-P)\le r-\lceil k/2\rceil-1$ and
$c(G)\le\bigl(q(G,P)+2^{-\lceil k/2\rceil-1}\bigr)2^{r}$. This is the lemma on cycles and a
marked set of [N3], with the integer part. The proof is the same count. I checked the two uses:
$|P|=2\lceil\log r\rceil$ gives the term $\frac1{2r}$ for the leading term, and five vertices
give $\frac25+\frac1{16}<\half$ in the last case. The copy [T] has eight vertices there, as [M];
both numbers are correct.

\subsection{Proposal 2: the window step}

The step uses three inequalities, and all three hold: $(t+\frac16)^{4}\le t^{4}+t^{3}$ for
$t\ge1$, because $\frac23+\frac16+\frac1{54}+\frac1{1296}<1$;
$6\mu^{3/4}+3k(6\mu^{1/4}+3)\le33\mu^{3/4}$ for $k\le\sqrt\mu$; and
$132\log^{2}\mu\le\mu^{1/4}$ for $\mu\ge2^{100}$, because $132\cdot10^{4}<2^{25}$ and the right
side grows faster. Step~1 of \M{Lemma~4.1} and the number $\Theta_\delta(k)$ do not change. This
is \T{Lemma~4.1, Step~2}.

\subsection{Proposal 3: a shortest path}\label{sec:path}

Let $N$ be a bridgeless $2$-pole of rank $b\ge2$, $A=w_0\cdots w_\ell$ a shortest path between
its ports, $J=N-V(A)$ and $\Delta=b-\ell$. Astra's statements, in the form of \T{Lemma~5.3}:
\begin{enumerate}[leftmargin=2.2em]
\item[(a)] $0\le\Delta\le\beta(J)$;
\item[(b)] at most $2\Delta$ layers between the ports are not thin, and at most $4\Delta$
vertices of $J$ lie outside the thin layers;
\item[(c)] every cycle of $J$ has at most $4\Delta+2$ vertices;
\item[(d)] $c(N)\le10\cdot2^{5\beta(J)+b/2}$.
\end{enumerate}
I checked (a), (b), (c) as written. In (d), the restriction of an even subgraph to $E(A)$ has
the kernel $\Z(J)$, so at most $2^{\beta(J)}$ cycles have the same edges on $A$; this is correct
also when an edge of $A$ is doubled.

\begin{idea}{the step that needs an exact reading}
Astra's notes say: if the layers at both ends of an edge $e_i$ of $A$ are thin and $e_i$ is
absent from the cycle, then the further edge of $w_{i-1}$ points forward and that of $w_i$
points backward, and ``positions satisfying this condition cannot be consecutive''. The
condition must be read as a property of the \emph{graph}: call $i$ free if the two further edges
have these directions. Then no two consecutive indices are free, an index that is not free
between thin layers has its edge in every cycle that passes it, and the count
$2^{\lceil(z-a-1)/2\rceil}$ is correct. Read as a property of the cycle (``no two consecutive
absent edges'') the count would be false: there are more than $2^{n/2}$ such patterns. The copy
[T] defines free indices.
\end{idea}

With (d), \M{Theorem~5.4} has a proof of five lines and the constant
$K_\varepsilon=20\cdot2^{5C_\varepsilon}$. \M{Proposition~5.3} keeps one conclusion,
$\beta(J)\le C_\varepsilon$, and Steps~1 to~3 of its proof. The ordinary layers, the count of
perfect matchings and the items (ii), (iii) are not necessary. Astra's notes keep the bound for
the layers that are not thin, because \M{Corollary~5.6} uses it. With [N2] that corollary is
not in the text, so (b) is used only inside the lemma.

In lines this proposal is neutral: Section~5 has 275 lines in [T] and 276 in the copy of [N3].
I put it in the copy for its content. The counts that hold for every pole are now in one lemma,
and the proposition with the forest estimate gives one number.

\subsection{Proposal 4: the gap theorem}

Astra's proof is that of [N3]: a bridgeless $2$-pole is a capped ladder if and only if $J$ is a
forest, and a cycle of $J$ gives a gap by \M{Lemma~5.1}. The two notes bound the length of the
cycle in different ways. [N3] takes it from Step~2 of \M{Proposition~5.3}. Astra takes it from
(c) and $\Delta\le\beta(J)\le C_{1/4}$, which gives $\gamma_0=2^{-4C_{1/4}-2}$. The copy [T]
has Astra's form (\T{Theorem~5.7}), because (c) is there in any case. Astra's warning is
correct: the old arithmetic with $3.99$ cannot use the new constant without change. [N3] changed
the arithmetic. The copy [T] does not need it (Section~\ref{sec:limits}).

\subsection{Proposal 5: caps for cuts of size two and three}

The proposal is correct. I checked the cap for two and for three ports, the inequalities
$z<4\eta$ and $\mathcal B>2^{d-2}-3\eta$, the step from $a_1+a_2\ge\frac23\sum_ja_j$ to
$a_3=o(1)$, the bounds $c(N)\le20\cdot2^{b/2}$ for a pole with a full pair and
$P_{ab}\le20\sqrt2\cdot2^{h/2}$ for a $3$-pole with two full pairs, and the constant: the excess
is at most $(15+400)\,2^{r/2}$.

I built it as an option of the script (\texttt{caps}). The copy has 1588 lines, 11 more than
[T]. The reason is that the proof of \M{Proposition~6.3} does not need full pairs. It uses
\M{Theorem~5.4} and one gadget, and it does not use the gap theorem or a cap. The proposal has
its merits: one method for all cuts, an absolute constant $420$, and no numbers $\frac1{1000}$,
$K_0$, $\kappa(h)$ in the paper. KNT can choose it; the proofs are in the option.

\subsection{Proposal 6: rigidity by limits}\label{sec:limits}

\emph{The three implications.} For the limit of the states of a sequence of $4$-poles whose
ranks tend to infinity, Astra states
\begin{gather*}
  a_{ik}+a_{jk}>1\ \Longrightarrow\ a_{ij}=0,\qquad\qquad
  a_{kl}>0\ \Longrightarrow\ a_{ij}+F-u_{ij|kl}\le1,\\
  a_{ik}+a_{il}+a_{jk}+a_{jl}>2\ \Longrightarrow\ u_{ik|jl}+u_{il|jk}=0 .
\end{gather*}
They need the gadgets of \M{Lemma~5.7} for every positive margin, as Astra says. So
\T{Lemma~5.9} is stated for every $\varepsilon$; its proof is that of [M].

\emph{The classification.} I checked the proof that a limit with $S+F\ge4$ is a cross or a star,
line by line, and a script enumerated all states on three grids (Test~3). There is no other
state.

\emph{Three full pairs.} Astra excludes them without \M{Lemma~6.7}: the pole without the two
other ports is a capped ladder, the two other ports cannot lie in its first end piece, and this
piece then has three edges to the rest of the graph. I checked the proof, and Test~2 examined
all positions of two ports in the capped ladders of weight $3$ to $6$.

\emph{What [T] takes.} The copy has the theorem for a cycle $C$ of length four (reduction~2 of
[N3]). There the count is
$c(G)\,2^{-k}=2^{-k}+z+(S+F-u_{12|34})+(S+F-u_{14|23})$, and one use of the cap gives
$S+F\ge4-\rho_0$. So [T] has: the cap with a free error $\eta$ (\T{Lemma~6.4}); the rigidity
lemma with the statement of [M] and Astra's proof by a limit (\T{Lemma~6.5}); and Astra's cut
of size three as Case~1 of \T{Theorem~6.6}. The second cap, the bound for $z$, the lemma on a
full star and the arithmetic of the star case are not necessary. Astra's step~6.5 (the full
pairs of the smaller side) is not necessary either, because the other side is $C$.

\emph{The constants.} In [N3] the gap constant forced $\varepsilon_0=\gamma_0/20$ in all of
Section~6. In [T] the gap theorem is used in the last step of the limit proof only. So
\M{Proposition~6.3} keeps its proof and its numbers $0.501$ and $1.002$.

\subsection{Proposal 7: the last section}

\emph{7.1, nesting.} The proof with the two cases $D\not\subseteq U_C$ and $D\subseteq U_C$ is
correct, and it is in [T]. \emph{Connected blocks:} for $j\ge i+2$ the separators between $C_i$
and $C_j$ lie in one component of $H-(C_i\cup C_j)$; I checked the argument with a path from $u$
to $p$. This removes the choice of a large component, and the factor $2\Lambda$ in the
thresholds. The proof of the one component has the same length as the choice of a component:
the option \texttt{blocks} gives 1578 lines.

\emph{7.2, two links.} The bound $W_2\le2\mu(e+NM)$ is correct, with $M$ a bound for the weight
of the regions that separate a neighbour $u$ of $A$ from $p$. It is the lemma of [M] in the
other direction, with no condition on $\mu$ and no number $K$. It is \T{Lemma~7.9}, and the
proof for heavy families is shorter with it. Test~4 compares the two sides for random families.

\emph{7.3 to 7.5.} I checked the packing bound $C_2d^{6}(s+1)2^{d}$ with $m=1$, the recursion
$|S_d|+1\le2^{2d}$ with $a(d)=\lceil C_0\log(Qd+2)\rceil$, and the end with $\mu\ge\frac53$ and
$\PP(Z>0)\ge\frac58-o(1)$. They are correct. The recursion from $d$ to a length of order
$\log d$ is the idea of [N1]. The two forms differ in one point: Astra deletes two nested sets
$S^{-}\subseteq S^{+}$, and [N1] deletes one set and counts the blocks that contain a short
cycle. The direct end is the inequality of the proposition on heavy families. I keep the form
of [N1], because [N2] is written for it: in [N2] the start of the recursion is girth five, not
a large girth $D_*$.

\subsection{Proposal 8 and the list of retained parts}

Proposal~8.1 is correct and has no use when there are no dyadic scales. Proposal~8.2 is correct;
Astra does not recommend it, and I did not build it.

Astra's notes keep \M{Lemma~5.5}, \M{Corollary~5.6} and \M{Section~7}, and they say that no
replacement was found. [N2] is such a replacement: a large piece of girth five has few forests,
so \M{Theorem~8.1} holds for girth five and the case of bounded girth is empty. Astra's notes do
not refer to [N2], so I take it that Astra has not audited it.

\section{The combined copy}\label{sec:copy}

\begin{center}\small
\begin{tabular}{@{}p{2.6cm}p{6.6cm}p{5.0cm}@{}}
\toprule
\textsf{\textbf{Section of [T]}} & \textsf{\textbf{Content}} & \textsf{\textbf{Origin of the changes}}\\
\midrule
3 Forest estimate & as [M]; Lemma~3.7 on a marked set; the leading term from it & [N3], Astra~1\\
4 Layers & Lemma~4.1 with the new Step~2 & Astra~2\\
5 Poles & Lemma~5.3 (a shortest path); Proposition~5.4 ($\beta(J)\le C$); Theorem~5.5 (few
cycles); Lemma~5.6 (ladders, three equivalent statements); Theorem~5.7 (gap); Corollary~5.8;
Lemma~5.9 (three gadgets, every margin) & Astra~3, 4; [N3]\\
6 Rings and cuts & rings as [M]; Proposition~6.3 as [M], also for loops; Lemma~6.4 (cap);
Lemma~6.5 (rigidity, by a limit); Theorem~6.6 (a cycle of length four) & [N3]; Astra~6\\
7 Girth five & lean form of [N2]; Lemma~7.7 (nesting) with the short proof; Lemma~7.9 (two
links) as an upper bound & [N1], [N2]; Astra~7.1, 7.2\\
8 Proof & three cases & [N3]\\
\bottomrule
\end{tabular}
\end{center}

\noindent Against [M], the copy does not have: Section~7 (bounded girth), \M{Lemma~5.5},
\M{Corollary~5.6}, \M{Lemma~5.9} (a third port), \M{Lemma~6.4} (gluing), \M{Lemma~6.7} (full
star), \M{Theorem~8.8} (raising the girth), the ordinary layers, and the edge recursion. It has
these new statements: the lemma on a marked set, the lemma on a shortest path, the two lemmas
of [N2] on full pairs and rigid sets at girth five with the lemma on pieces, and the proposition
on heavy families.

\subsection{Measured}

\begin{center}\small
\begin{tabular}{@{}p{5.2cm}rrrrr@{}}
\toprule
 & \textsf{\textbf{[M]}} & \textsf{\textbf{[N3] copy}} & \textsf{\textbf{[T]}} & \textsf{\textbf{[T], caps}} & \textsf{\textbf{[T], blocks}}\\
\midrule
Pages & 42 & 41 & 41 & 41 & 41\\
Typeset lines & 1624 & 1587 & 1577 & 1588 & 1578\\
Numbered statements & 40 & 39 & 39 & 39 & 39\\
Cases of the main theorem & 5 & 3 & 3 & 3 & 3\\
Section 4 (layers), lines & 72 & 73 & 68 & 68 & 68\\
Section 5 (poles), lines & 339 & 276 & 275 & 278 & 275\\
Section 6 (rings and cuts), lines & 275 & 258 & 252 & 260 & 252\\
Bounded girth and last section, lines & 323 & 358 & 355 & 353 & 356\\
Proof of the main theorem, lines & 30 & 20 & 20 & 20 & 20\\
\bottomrule
\end{tabular}
\end{center}

\noindent The numbers of statements are equal because one lemma comes (a shortest path) and one
goes (a full star). The gain of [T] against the copy of [N3] is ten lines. Its larger gain is in
the content: no numbers that depend on the gap constant, no lemma on a full star, and the
counts in a pole separated from the forest estimate.

\section{The text and the formal proof}\label{sec:formal}

Nothing of Astra's notes or of my notes is formalised. In [M] the text and the Lean proof differ
in four places. In [T] they differ in these places also; Section~1.4 of [T] says so.
\begin{itemize}[leftmargin=1.2em]
\item The leading term: the formal proof has the edge recursion.
\item \T{Theorem~5.5}: the formal proof counts perfect matchings between ordinary layers.
\item \T{Theorem~5.7}: the formal proof is an induction with a third port, and it gives
$\frac1{16}$.
\item \T{Lemma~6.5} and \T{Theorem~6.6}: the formal proof treats a cyclic cut of size four with
two sides; it has the arithmetic with fixed numbers and the bound for a full star.
\item The last case: the formal proof has the case of bounded girth and two girth scales.
\end{itemize}
The formal statement of the main theorem does not depend on any of this.

\section{Tests}\label{sec:tests}

The script \texttt{checks/astra\_audit\_check.py} has six tests. Exit code $0$ means that all
passed. With the arguments \texttt{3000 2} it prints:

\begin{center}\small
\begin{tabular}{@{}p{0.8cm}p{13.7cm}@{}}
\toprule
\textsf{\textbf{Test}} & \textsf{\textbf{Content and result}}\\
\midrule
1 & \T{Lemma~5.3} for 3000 random bridgeless $2$-poles of rank $2$ to $7$ and all 4157 shortest
paths: (a), (b), (c); at most $2^{\beta(J)}$ cycles with the same edges on $A$; no two
consecutive free indices; the edge of an index that is not free between thin layers is in the
cycle; the bound of~(d). The largest value of $c(N)\,2^{-\beta(J)-4\Delta-\ell/2}$ was $4.125$;
the lemma allows $10$.\\
2 & 9354 capped ladders of weight $3$ to $6$ with two more ports on their edges. In 282 of them
both pairs at the cap are full. In each of these, no port is in the first end piece, and two
edges leave it.\\
3 & All states with entries in the grids of step $\frac14$, $\frac16$, $\frac18$ that satisfy
the three implications, $F\le1$ and $S+F\ge4$: 361 states, 9 crosses and 352 stars.\\
4 & $W_2\le2\mu(e+NM)$ for 300 random families of regions with weights in cubic graphs (5752
ordered pairs with two links). The largest ratio of the two sides was $0.263$.\\
5 & The two inequalities of the window step, for $t\in[1,101]$ and $\log\mu\in[100,20000]$.\\
6 & The caps of poles with two and three ports: the count of the cycles through the cap, 150
poles.\\
\bottomrule
\end{tabular}
\end{center}

\noindent Tests 1, 2, 4, 6 are complete counts on small graphs. Test~3 is a finite part of a
statement about real numbers. The rigidity lemma and the theorem on a cycle of length four
contain large constants, and no computation tests them.

\section{Open points}\label{sec:open}

\begin{itemize}[leftmargin=1.2em]
\item \emph{The base.} The copy [T] uses [N2] (no case of bounded girth) and the cycle of length
four of [N3], which needs [N2]. Astra's notes are written for [M] with its Section~7. The
proposals 2, 3, 4 and 7.2 do not depend on this choice. Proposal~6 then needs its form with two
sides (Astra's steps 6.3 to 6.5, which I checked on paper), and \M{Corollary~5.6} needs (b) of
the lemma on a shortest path. I did not build this copy.
\item \emph{Readers.} [N2] has had one reader. Its two lemmas at girth five carry the removal of
a whole case, and they have very large constants. A second reader for them has more value than
for any other part.
\item \emph{The formal proof.} A formal proof of [T] is new work in Sections~5 and~6 and in the
last section. The existing Lean proof stays valid for the theorem.
\item \emph{The options.} \texttt{caps} and \texttt{blocks} are complete texts. Each is a choice
of KNT; I recommend neither, by the measured lines.
\end{itemize}

\section{Files}\label{sec:files}

\begin{center}\small
\begin{tabular}{@{}p{6.1cm}p{8.4cm}@{}}
\toprule
\textsf{\textbf{File}} & \textsf{\textbf{Content}}\\
\midrule
\texttt{research/astra-audit.tex}, \texttt{.pdf} & this note\\
\texttt{research/apply\_combined.py} & writes the combined form into a copy; options
\texttt{caps} and \texttt{blocks}; every replacement must match exactly once\\
\texttt{research/trial-combined.pdf} & the copy [T], 41~pages\\
\texttt{checks/astra\_audit\_check.py} & the six tests of Section~\ref{sec:tests}\\
\bottomrule
\end{tabular}
\end{center}

\noindent To build [T], copy the folders \texttt{exact-maximum} and \texttt{style} to a new
folder. On the copy of \texttt{exact-maximum} run the three scripts in this order:
\begin{center}
\texttt{python3 apply\_girth\_five.py <copy> lean}\\
\texttt{python3 apply\_fewer\_cases.py <copy> keepgap}\\
\texttt{python3 apply\_combined.py <copy>}
\end{center}
Then build \texttt{exact-maximum.tex} with \texttt{tectonic}.

\begin{thebibliography}{9}\small
\bibitem[M]{M} KNT, Astra and Fable~5.2, \emph{The maximum number of cycles in a graph of large
rank}, draft of 12 November 2026, 42~pages, \texttt{exact-maximum/exact-maximum.pdf}.
\bibitem[T]{T} The combined copy of this note, \texttt{research/trial-combined.pdf}.
\bibitem[N1]{N1} Fable~5.2, \emph{The large-girth theorem with one family of short cycles},
\texttt{research/one-scale.pdf}.
\bibitem[N2]{N2} Fable~5.2, \emph{Pieces at girth five: the case of bounded girth is not
necessary}, \texttt{research/girth-five.pdf}.
\bibitem[N3]{N3} Fable~5.2, \emph{Three cases: four reductions in the cut analysis and the final
count}, \texttt{research/fewer-cases.pdf}.
\end{thebibliography}

\end{document}
